\subsection{Weighted curvature of a completed passage}

To use an action calculation in the physical input, we must include the
curvature contributed by each section map. The following identity keeps
these contributions separate and allows a positive weight adapted to
the coefficient chart.

\begin{lemma}[Weighted curvature of a completed passage]
\label{neg:lem-completed-curvature}
Let \(u,v>0\) be radial coordinates, with \(P_\pm,B_\pm>0\)
on the two strict through intervals, and suppose
\[
 v_u=\frac{uB_+(u)P_-(v)}{vB_-(v)P_+(u)}(1+\varrho)>0,
 \quad S(-u^2/2)=-v^2/2,\quad
 \Pi=\phi_-\circ S\circ\phi_+^{-1},
\]
Here \(P_\pm,B_\pm\) are \(C^2\), \(\varrho=\varrho(u,v)\)
is \(C^2\) on the product of the two radial intervals, and
\(\phi_\pm\) are the separately normalized actual \(C^3\) section
maps with positive derivatives. Put
\(\mathcal D_+=P_+\partial_u/(uB_+)\),
\(\mathcal D_-=P_-\partial_v/(vB_-)\), and \(R_B=B_+/B_-\).
For any smooth positive weight \(W=W(u,v)\), one has exactly
\begin{align}
 \phi_+'\frac{P_+}{W}\frac{\Pi''}{\Pi'}
 ={}&\frac{P_+' /u-R_B(1+\varrho)P_-'/v}{W}
 -\frac{P_+}{uW}\frac d{du}\log R_B \notag\\
 &-\frac{B_+}{W}
       \{\mathcal D_++(1+\varrho)\mathcal D_-\}\log(1+\varrho)
       \notag\\
 &-\frac{B_+(1+\varrho)}{W}\mathcal D_-\log\phi_-'
 +\frac{P_+}{u^2W}(u\partial_u)\log\phi_+'.
 \label{deg:weighted-exact-curvature-chain}
\end{align}
Primes on each \(\phi_\pm\) use its own quadratic label; primes on
\(P_\pm\) use the displayed radial variable. The derivative \(d/du\)
follows the matched graph. In particular both section terms remain in
the curvature of the actual affine input.
\end{lemma}
\begin{proof}
One has \(S'=R_BP_-(1+\varrho)/P_+\). Logarithmically differentiate
this identity using \(d/db_+=-u^{-1}d/du\), and substitute in
\[
 \phi_+'\frac{\Pi''}{\Pi'}=
 \frac{S''}{S'}+S'\frac{\phi_-''}{\phi_-'}-
                         \frac{\phi_+''}{\phi_+'}.
\]
Multiplication by \(P_+/W\) gives the formula. Its coefficients need
not extend without weights to a stopped endpoint.
\end{proof}

